% =============================================================================
%  Section: Summary and Conclusions
%  H1 Charged-Hadron SIDIS High-Q^2 Analysis Note
% =============================================================================
\section{Summary and Conclusions}
\label{sec:summary}

This note documents the inclusive charged-hadron SIDIS analysis at high
$\Qsq$ with H1.  The analysis uses the full four-period HERA-II data sample with
an integrated luminosity of $351~\pbinv$ and measures charge-separated inclusive
hadrons without PID splitting.  The truth-level charged-hadron convention for
response and acceptance work is limited to $\pi^{\pm}$, $K^{\pm}$, and $p/\bar p$.

The current analysis release has a validated 490-row published-bin contract,
with 245 $\hplus$ and 245 $\hminus$ bins, $N_{B,G}=245$ and $N_{B,R}=490$ per
charge setup, response shape $(491,246)$ per charge, data-vector length 980,
and a nominal four-iteration iterative Bayesian unfolding.  The selection
includes the $\alpha\ge0.2~\mathrm{rad}$ laboratory-frame scattered-electron
cone on the reconstructed legs, the unfolding prior is the response matrix's
own generator-level distribution, and the electron, V0, and low-$\Qsq$
background templates are subtracted from the data vector before unfolding.
Since the 2026-10-20 OPT-A baseline promotion, the response truth contract is
harmonized to the radiative correction
(Sec.~\ref{sec:unfolding:truth-contract}), and the OPT-A chain is the nominal
central-value release (Sec.~\ref{sec:results:z005-release}).
Detector, unfolding, tracking-efficiency,
luminosity, and radiative systematics have been propagated into the total
summary.  A
RooUnfoldBayes cross-check reproduces the in-house IBU central yields to
floating-point precision, with a maximum relative difference of
$2.6\times10^{-15}$ across 980 rows.

The radiative correction is applied to all released cross sections, which are
therefore Born-level, and its uncertainty is carried as two systematic
components; the application of the correction at the multiplicity level remains
open.  The principal remaining limitations are this multiplicity-level mapping
and the absence of a corrected inclusive DIS denominator.
The current release should therefore be treated as a validated working release,
not yet as a final physics multiplicity release.  The quoted total systematic
budget of the frozen 2026-10-19 edition combines twenty active uncertainty
terms on the background-subtracted basis, with median totals of
$7.320288331169215\%$ ($\hplus$) and $7.637124635311324\%$ ($\hminus$); the
frozen editions were computed on the superseded central values and are
superseded as a quoting basis, the budget is being regenerated on the OPT-A
nominal, and no per-bin uncertainty or total is final until the regenerated
budget passes review (Sec.~\ref{sec:syst:budget-regeneration}); its
principal background-related term is a fully correlated $\pm100\%$
normalization envelope on the subtracted electron and V0 content.  The five
detector-shift components are full-chain single-sided variations; the HFS legs
are produced in the plus direction only and their sign symmetry is assumed,
while the tracking-efficiency component remains a data-side conservative
shift.  The electron, V0, and low-$\Qsq$ backgrounds are subtracted from the
data vector before unfolding, while the detector-interaction removable class is
removed response-internally and must not be subtracted again externally.  The
applied $C_{\mathrm{rad}}$ is
mixed-beam only, and the release is not publication-grade: the exposure of the
simulated radiative samples is reconciled file by file, but the generator-log
cross-check is permanently blocked externally, and the correction tables
remain provisional.  The lab-$p_T$, $|\eta|$, and DCA variations provide
robustness checks rather than central results.

The low-$\Qsq$ charged-SIDIS companion analysis has a separate preliminary
v3p3p1 TUnfold publication-track result and open systematic follow-ups.  It is
kept separate from the high-$\Qsq$ analysis and is summarized here only to make
the charged-SIDIS status across the two regimes explicit.
